% ! TeX root = ../main.tex

% Background: Volatile function tail latency example
\newcommand{\figBgVolatile}{%
  \begin{figure}[t]
    \centering
    \includegraphics[width=\columnwidth]{lesson_volfunc_zoom.pdf}
    % \caption{Simplified view of telemetry from a real BSP app with a volatile function followed by a global collective.
    %   Tail latency in the volatile function gets misattributed as elevated collective time in aggregate
    %   telemetry (left), but not in fine-grained telemetry (right).
    % }
    % \caption{Telemetry from an AMR code where a volatile \mpiw{} precedes \texttt{MPI\_Allgather}.
    %   Aggregate telemetry (left) misattributes straggler-induced \mpiw{} as elevated collective cost.
    %   Fine-grained telemetry (right) makes the cause obvious but requires joining all ranks' telemetry.
    % }
    \caption{Reproduced performance pathology from an AMR code~\cite{\citeJainAmr}.
      A volatile \mpiw{} preceding \texttt{MPI\_Allgather} is misattributed as elevated collective cost
      in aggregate telemetry (left).
      Rank-level telemetry for a single timestep (right) reveals the true cause but requires joining all
      ranks' data.
    }
    \label{fig:orca:bg-volatile}
  \end{figure}%
}

% Design: ORCA Architecture
\newcommand{\figDesignArch}{%
  \begin{figure}[t]
    \centering
    \includegraphics[width=0.48\textwidth, trim=0.2cm 0cm 0.2cm 0cm, clip]{thesis-orca-arch.pdf}
    % \caption{ORCA architecture.
    %   A three-tier TBON with application ranks (\tiermpi{}) at the leaves, aggregator nodes (\tieragg{})
    %   at intermediate levels, and a controller (\tierctl{}) at the root.
    %   Data flows upward as columnar Arrow streams, transformed en-route by OrcaFlow operators, and
    %   persisted via Parquet or DuckDB sinks.
    % }
    \caption{ORCA architecture.
      A three-tier TBON with application ranks (\tiermpi{}) at the leaves, intermediate aggregators
      (\tieragg{}), and a controller (\tierctl{}) at the root.
      Telemetry originates as columnar timestep dataframes, is transformed by SQL plan operators at each
      tier, and persisted via DuckDB and Parquet sinks --- enabling workflows from real-time feedback via
      lightweight metrics to zero-ETL post-hoc analytics.
    }
    \label{fig:orca:des-arch}
  \end{figure}%
}

% Design: ORCA Flow Plan
\newcommand{\figDesignPlan}{%
  \begin{figure}[t]
    \centering
    \includegraphics[width=0.24\textwidth, trim=0.35cm 0cm 0.2cm 0.4cm, clip]{plan.pdf}%
    \includegraphics[width=0.24\textwidth, trim=0.3cm 0cm 0.3cm 0.4cm, clip]{plan2.pdf}
    % \caption{OrcaFlow compilation. (Left) An OrcaFlow definition specifying SQL transforms over
    %   telemetry streams. (Right) The compiled tier plans showing operator placement across
    %   \tiermpi{}, \tieragg{}, and \tierctl{}, with filters pushed to leaves and aggregation
    %   at higher tiers.}
    \caption{OrcaFlow compilation showing flows and their compiled tier plans.
      Flows chain SQL transforms over telemetry streams en route to configurable sinks.
      Both operators and sinks are placed automatically.
      (Left) A selective filter pushed to \tiermpi{}, discarding most data at source.
      (Right) Aggregations also reduce data volume --- pushed-down partial aggregations emit fixed-size intermediates regardless of input volume.
    }
    \label{fig:orca:des-plan}
  \end{figure}%
}

% Design: \twopc{} Protocol
\newcommand{\figDesignTwopc}{%
  \begin{figure}[t]
    \centering
    \includegraphics[width=0.48\textwidth, trim=0.2cm 0cm 0.2cm 0cm, clip]{ts2pc_paper.pdf}
    \caption{Timestep-linked two-phase commit (\twopc{}).
      Commands are proposed for a future timestep $T_{cmd} = T_{mpi} + \Delta$.
      (A) Normal case: ranks agree and apply the command before executing $T_{cmd}$.
      (B) Abort: ranks have passed $T_{cmd}$ and vote to abort, command is retried with a
      larger $\Delta$. (C) Stall: ranks have prepared but not received commit/abort.
      Execution blocks until resolution to preserve correctness.
    }
    \label{fig:orca:des-twopc}
  \end{figure}%
}

% Eval: Runtime, trace size, query suite (3 subfigures)
\newcommand{\figEvalBannerAll}{%
  \begin{figure*}[t]
    \begin{subfigure}[t]{\textwidth}
      \centering
      \includegraphics{amr_tracer_runtimes.pdf}
      \caption{Runtime overhead.}
      \label{fig:orca:eval-banner-runtime}
    \end{subfigure}
    \begin{subfigure}[t]{0.49\textwidth}
      \centering
      \includegraphics[width=\textwidth]{amr_tracesizes_line.pdf}
      \caption{Storage efficiency.}
      \label{fig:orca:eval-banner-size}
    \end{subfigure}
    \hfill
    \begin{subfigure}[t]{0.49\textwidth}
      \centering
      \includegraphics[width=\textwidth]{query_suite_abs.pdf}
      \caption{Query performance.}
      \label{fig:orca:eval-banner-query}
    \end{subfigure}
    \caption{ORCA substrate evaluation at 512--4096 ranks. (Top) Runtime overhead: ORCA maintains
      0.8--2\% while traditional tracers degrade to 56--81\% at scale. (Bottom-left) Storage:
      Parquet is 2--20$\times$ smaller than alternatives. (Bottom-right) Query latency: ORCA is
      3--4 orders of magnitude faster due to columnar format and partition pruning.}
    \label{fig:orca:eval-banner}
  \end{figure*}%
}

% Eval: Runtime overhead, two-column
\newcommand{\figEvalOverhead}{%
  \begin{figure*}[t]
    \centering
    \includegraphics[width=\textwidth]{amr_tracer_runtimes.pdf}
    % \caption{Runtime overhead comparison at 512--4096 ranks.
    %   ORCA maintains 0.8--2\% overhead across all scales and profiles, while traditional tracers degrade
    %   from 18--20\% at 512 ranks to 56--81\% at 4096 ranks due to uncoordinated buffer flushes.
    % }
    \caption{Runtime overhead for a detailed tracing profile on an AMR code at 512--4096 ranks.
      An additional minimal profile (collective-only) for ORCA shows always-on overhead.
      ORCA maintains 0.8--2\% overhead across both profiles, while TAU and \dftracer{} degrade from
      18--20\% at 512 ranks to 56--81\% at 4096 ranks.
      \emph{*extrapolated from 200 timesteps due to stability issues.}
    }
    \label{fig:orca:eval-runtime}
  \end{figure*}%
}

% Eval: Trace size only, one-column
\newcommand{\figEvalSizeOnly}{%
  \begin{figure}[t]
    \centering
    \includegraphics[width=\columnwidth]{amr_tracesizes_line.pdf}
    \caption{Trace size (left) and per-record size (right) comparison for detailed profile.
      ORCA traces are 3.6--44$\times$ smaller than other formats.
      Per-record size isolates format efficiency from record amplification --- other tracers emit
      2.5--3.3$\times$ more records than ORCA for the same collection profile.
    }
    \label{fig:orca:eval-tracesize}
  \end{figure}%
}

% Eval: Query suite only, one-column
\newcommand{\figEvalQueryOnly}{%
  \begin{figure}[t]
    \centering
    \includegraphics[width=\columnwidth]{query_suite_abs.pdf}
    \caption{Post-mortem query latencies on 20 timesteps.
      ORCA queries execute in O(100 ms), 3--4 orders of magnitude faster than
      \dftracer{} and Caliper.
      Performance for plaintext formats is dominated by parsing, not query complexity.
      TAU and \scorep{} are omitted as OTF2 is not a queryable format.
    }
    \label{fig:orca:eval-query}
  \end{figure}%
}

% Eval: Architecture comparison (OFI/TCP + DFTracer compression)
\newcommand{\figEvalArch}{%
  \begin{figure}[t]
    \centering
    \begin{subfigure}[t]{0.65\columnwidth}
      \centering
      \includegraphics{dftracer_comp.pdf}
      \caption{Compression impact.}\label{fig:orca:eval-arch-dft}
    \end{subfigure}
    \hfill
    \begin{subfigure}[t]{0.33\columnwidth}
      \centering
      \includegraphics{ofi_tcp_tracer_runtimes.pdf}
      \caption{Transport impact.}\label{fig:orca:eval-arch-tcp}
    \end{subfigure}
    % \caption{Architectural analysis at 512 ranks. (Left) Enabling compression for \dftracer{}
    %   reduces trace size but increases runtime overhead due to CPU contention; ORCA's columnar
    %   format achieves comparable compression without this tradeoff. (Right) RDMA transport
    %   achieves lower overhead than TCP by bypassing the kernel.}
    \caption{(Left) Compression reduces trace size but induces severe jitter on simulation ranks.% --- \dftracer{} with gzip degrades to 504\% overhead.
      ORCA avoids this by offloading compression to dedicated aggregators.
      (Right) TCP transport increases ORCA overhead from 1--2\% to 45--55\%.
      Arrow-RDMA co-design is essential for low overhead.
    }
    \label{fig:orca:eval-arch}
  \end{figure}%
}

% Eval: Online vs offline
\newcommand{\figEvalOrcaNative}{%
  \begin{figure}[t]
    \centering
    \includegraphics[width=\columnwidth]{ontv_combined.pdf}
    % \caption{In-situ analytics with operator pushdown.
    %   Runtime overhead remains below 2\% across all flows despite executing pushed-down operators on MPI
    %   nodes.
    %   Pushdown reduces transmitted data volume by 98.5--99.999\%, enabling continuous observability that
    %   would be impractical with capture-everything tracing.
    % }
    \caption{Runtime overhead (top) and MPI tier data reduction (bottom) for in-situ analytics with operator pushdown.
      Flows range from monitoring communication outliers to compute load balance.
      Pushdown discards 98.5--99.999\% of data at sub-2\% overhead, despite operators executing on MPI
      ranks.
    }
    \label{fig:orca:eval-ontv}
  \end{figure}%
}

% Eval: Verbs vs TCP
\newcommand{\figEvalOrcaTcp}{%
  \begin{figure}[t]
    \centering
    \includegraphics[width=\columnwidth]{ofi_tcp_tracer_runtimes.pdf}
    \caption{Transport comparison: RDMA verbs vs TCP.
      RDMA achieves lower and more consistent overhead by bypassing the kernel and avoiding jitter from
      socket operations.
    }
    \label{fig:orca:eval-tcp}
  \end{figure}%
}

% Eval: Agentic debugging timeline
\newcommand{\figEvalAgentic}{%
  \begin{figure}[t]
    \centering
    \includegraphics[width=\columnwidth]{agentic_timeline.pdf}
    % \caption{Agentic debugging session. (Top) User and LLM interactions over time, with numbered
    %   flow installations. (Bottom) Per-timestep event volume showing how the agent progressively
    %   narrows collection scope---from monitoring collectives to filtering Kokkos events to
    %   capturing targeted traces---localizing the anomaly from \Cref{fig:orca:bg-volatile} in minutes.}
    \caption{Visual transcript of an agentic debugging session demonstrating closed-loop capabilities.
      An LLM agent is tasked with localizing the anomaly from \Cref{fig:orca:bg-volatile} using high-level
      metrics, progressively collecting targeted telemetry to identify the source.
      The agent operates largely autonomously, producing the complete call path in 27 minutes.
      Reverting to baseline monitoring reduces data volume by 144$\times$.
    }
    \label{fig:orca:eval-agentic}
  \end{figure}%
}

% Eval: In-situ analytics flows table
\newcommand{\tablesink}[1]{\textcolor{gray!80}{$\rightarrow$ \emph{#1}}}
\newcommand{\tablevar}[1]{\texttt{\textit{\textcolor[HTML]{4D6B5A}{#1}}}}
\newcommand{\tabEvalFlows}{%
  \begin{table}[t]
    \centering
    \footnotesize
    \begin{tabular}
      {@{}p{2.3cm}p{5.6cm}@{}}
      \toprule
      \textbf{Flow}         & \textbf{Query}                                                            \\
      \midrule
      Outlier Wait~Counts   &
      \texttt{SELECT COUNT(*) FROM mpi\_wait WHERE~duration > 1ms} \tablesink{DuckDB}                   \\
      \addlinespace
      Outlier Wait~Traces   &
      \texttt{SELECT * FROM mpi\_wait WHERE~duration~>~1\,ms} \tablesink{Parquet}                       \\
      \addlinespace
      Load \mbox{Imbalance} &
      \texttt{SELECT rank, SUM(duration) FROM kokkos WHERE name IN \textcolor{gray}{<compute kernels>}} \\
                            & \texttt{GROUP BY rank} \tablesink{\tablevar{stats}} \tablesink{Parquet}   \\
      \addlinespace[0.2em]
                            & \texttt{PERCENTILES(\tablevar{stats}, 50/90/99\textsuperscript{th})} \tablesink{DuckDB}     \\
      \bottomrule
    \end{tabular}
    \caption{In-situ workflows evaluated in \S\ref{sec:orca:eval-insitu}, ranging from lightweight metrics and selective tracing on the communication path to aggregates over compute kernels.}\label{tab:ontv-flows}
  \end{table}%
}